\section{Conclusion}
\label{sec:conclusion}

We replaced the winner-takes-all checkpoint selection that ends most vision
pipelines with a rule that costs nothing at training or inference time: rank
the checkpoints of a completed run by the validation criterion the pipeline
already computes, keep the top $k$, average them. Because the rule reads only
one scalar per checkpoint, oriented by a single sign, the same procedure
applies to classification, detection and segmentation unchanged.

Across 34 model--dataset configurations with five seeds each, under a protocol
in which the baseline and every $k$ share a training trajectory, $k{=}5$
improves the primary metric in all 34: classification accuracy rises
81.73\%$\to$83.76\% and macro-F1 81.65\%$\to$83.86\% in equally weighted mean,
detection and segmentation fitness gain 5.8--6.9\% relative, and across-seed
standard deviation drops by roughly a fifth. The effect of $k$ is not uniformly
monotonic per configuration, so the finding is not that larger $k$ is always
better, but that a small criterion-ranked set is a better default than a single
winner. We also showed the rule is not last-$k$ averaging in disguise: the
selected checkpoints are usually non-contiguous and rarely near the end of
training. The matched comparison that observation invites---last-$k$, EMA,
post-hoc SWA and a BN-recalibrated baseline on the same trajectories---is the
natural next step.
